\section{Introduction}
\label{sec:vof}

\hypertarget{cs_vof}{}

The Volume of Fluid (VOF) model allows for the simulation of immiscible
multi-phase flows with a resolved interface. It belongs to the class of
interface-capturing approaches in which a scalar indicator function is
transported.

In \CS, the VOF formulation handles two phases and relies on the introduction
of an additional transport equation for the void fraction of one of the phases.
Mass and momentum equations remain unchanged, the velocity corresponding to the
velocity of the mixture. The fluid properties of density and viscosity are
reconstructed from the local void fraction, using constant values for each
phase. Surface tension can be enabled through a Continuum Surface Force
approach, and interface compression can be activated.

Dedicated compressive numerical schemes with boundedness properties are used
for the void fraction transport. The transport equation is solved in
conservative form so that the total amount of each phase is preserved at the
discrete level.

The present approach is especially suited to free-surface flows, wave dynamics,
sloshing, dam-break problems, liquid-gas interactions and cavitating flows.

%-------------------------------------------------------------------------------

\section{System equations}

\subsection{Volume fraction transport equation}

Let $\alpha$ denote the void fraction of the dispersed phase
(or equivalently the void fraction in cavitation applications).
%alpha=0 by default initialization, phase 1 to 1000 by default
%rho_0 taken to max (rho_1, rho_2)

The local phase indicator satisfies:

\begin{equation}
\frac{\partial \alpha}{\partial t}+\divs (\alpha \vect{u})=S_\alpha ,
\label{eq:vof:transport}
\end{equation}

where:

\begin{itemize}
\item $\alpha = 0$ corresponds to pure fluid 1,
\item $\alpha = 1$ corresponds to pure fluid 2,
\item $0 < \alpha < 1$ identifies cells crossed by the interface,
\item $\vect{u}$ is the mixture velocity,
\item $S_\alpha$ denotes optional source terms.
\end{itemize}

An additional drift contribution may be introduced:

\begin{equation}
\frac{\partial \alpha}{\partial t}
+\divs \cdot (\alpha \vect{u})
+\divs \cdot (\alpha (1-\alpha)\vect{u}_r)
=S_\alpha ,
\label{eq:vof:drift}
\end{equation}

where $\vect{u}_r$ denotes a relative velocity between phases.
This formulation is based on the drift-flux approach implemented in \CS
following the work of \cite{deshpande:2012}.
It can be either interpreted as :
\begin{itemize}
  \item A compressive contribution to preserve a sharper transition between
    phases as pure advection tends to diffuse the interface because of
    numerical dissipation, despite of compressive
    numerical schemes used.
  \item A specific closure to model two-phase behaviours at scales that are not
    represented by the mesh.
\end{itemize}
The additional flux is active only close to the interface since it is
proportional to $\alpha(1-\alpha)$.

%-------------------------------------------------------------------------------

\subsection{Mixture properties}

The VOF model adopts a one-fluid formulation.
Only one velocity field is solved on the computational domain.

The local density is reconstructed from the void fraction using a
linear interpolation:

\begin{equation}
\rho=\alpha \rho_1+(1-\alpha)\rho_2 ,
\label{eq:vof:rho_mix}
\end{equation}

with $\rho_1$ and $\rho_2$ the reference viscosities of the phases 1 and 2.

Similarly, the dynamic viscosity is obtained as

\begin{equation}
\mu=\alpha \mu_1+(1-\alpha)\mu_2 ,
\label{eq:vof:mu_mix}
\end{equation}

with $\mu_1$ and $\mu_2$ the reference viscosities of the phases 1 and 2.

%-------------------------------------------------------------------------------

\subsection{Mixture mass conservation}

The mixture density varies locally through the void fraction field.
The corresponding conservation equation reads

\begin{equation}
\frac{\partial \rho}{\partial t}+\divs \cdot (\rho \vect{u})=0.
\label{eq:vof:mass}
\end{equation}

The mass fluxes used by the momentum solver are updated consistently
with the reconstructed density field so that phase conservation and
momentum conservation remain coupled.

%-------------------------------------------------------------------------------

\subsection{Momentum equation}

The velocity field is governed by a single Navier--Stokes equation
written for the mixture:

\begin{equation}
\frac{\partial (\rho \vect{u})}{\partial t}
  +\divv (\rho \vect{u}\otimes\vect{u})
=-\grad p
+\divv\left[\mu\left(\gradt\vect{u}+\gradt\vect{u}^{T}\right)\right]
+\rho \vect{g}+\vect{F}_{\sigma},
\label{eq:vof:momentum}
\end{equation}

where:

\begin{itemize}
\item $p$ is pressure,
\item $\vect{g}$ is gravitational acceleration,
\item $\vect{F}_{\sigma}$ is the surface tension force.
\end{itemize}

The momentum equation is therefore solved over the whole domain,
independently of the local phase, while material properties vary through
the void fraction field.

Surface tension effects are modelled through the Continuum Surface Force (CSF)
formulation proposed by Brackbill, Kothe and Zemach. It spreads the capillary
force over a finite thickness around the interface.

The interface normal is computed from the void fraction gradient:

\begin{equation}
\vect{n}=\frac{\grad \alpha}{\left| \grad \alpha \right|}.
\label{eq:vof:normal}
\end{equation}

The interface curvature is obtained from

\begin{equation}
\kappa=-\divs \vect{n}.
\label{eq:vof:curvature}
\end{equation}

The volumetric surface tension force entering the momentum equation is

\begin{equation}
\vect{F}_{\sigma}=\sigma \kappa \grad \alpha,
\label{eq:vof:csf}
\end{equation}

where $\sigma$ is the surface tension coefficient.
